\begin{abstract}
RLHF and DPO take different optimization paths to align language models with human preferences: RLHF trains an explicit reward model, while DPO directly optimizes from preference pairs. We test whether these mechanistic differences produce distinct alignment signatures under controlled conditions using Llama-2-7B and Anthropic's HH-RLHF dataset.

Our five-hypothesis experiment validates the mechanistic premise: RLHF produces smooth, continuous reward predictions (range 0.83 units), while DPO preserves sharper preference boundaries (sharpness ratio 1.65, 65\% higher margin variance). However, these training dynamics do not translate to distinct downstream effects at 7B scale with LoRA fine-tuning. Models trained with different methods do not cluster into separate behavioral attractors---cross-method similarity actually exceeds within-method similarity (clustering gap = $-0.016$). Standard alignment benchmarks (TruthfulQA, HHH) show similar performance profiles (max $|d|$ = 0.194, profile correlation = 0.978).

We term this finding ``different dynamics, similar destinations'': verified mechanistic differences do not reliably manifest as distinct alignment signatures on existing evaluation infrastructure. This has practical implications for method selection and suggests limitations in current benchmark sensitivity.
\end{abstract}
